% ===========================================================================
% ΟΔΗΓΟΣ ΔΙΑΓΡΑΜΜΑΤΩΝ ΚΑΙ ΣΧΗΜΑΤΩΝ — ΚΟΙΝΟΣ ΣΕ PRINT ΚΑΙ SCREEN
% ===========================================================================
% Οι ίδιες γεωμετρίες και τα ίδια δεδομένα παράγουν και τις δύο εκδόσεις.
% Τα styles δηλώνουν ρόλο μέσα στο συγκεκριμένο σχήμα, όχι φυσικό μέγεθος.
\setcounter{chapter}{1}
\markboth{Οδηγός σχημάτων}{Οδηγός σχημάτων}
\thispagestyle{plain}
\begin{center}
{\BookHeadingFont\color{BookPrimary}\fontsize{17}{22}\selectfont Οδηγός διαγραμμάτων και σχημάτων}\par
\smallskip
{\color{BookInk}\small Μαύρο για αναφορά · δεύτερο μελάνι για έμφαση ανά σχήμα}
\end{center}
\BookSection{Α. Κοινές αρχές}
\begin{itemize}
\item \textbf{Ρόλοι:} άξονες, αντικείμενα και γράμματα μαύρα. Το κύριο στοιχείο κάθε σχήματος παίρνει το δεύτερο μελάνι.
\item \textbf{Διάκριση σειρών:} συνεχής και διακεκομμένη γραμμή, μαζί με ετικέτες· η πληροφορία παραμένει σαφής και χωρίς χρώμα.
\item \textbf{Γράμματα:} ετικέτες 9 pt, αριθμοί διαβαθμίσεων 8 pt. Μαθηματικά σύμβολα πλάγια και μονάδες SI όρθιες, π.χ. $x\,(\mathrm{m})$.
\item \textbf{Γεωμετρία:} άξονες και θετικές φορές δηλώνονται ανά πρόβλημα. Δεν επιβάλλεται ίδια κατακόρυφη φορά σε όλα τα θέματα.
\item \textbf{Τελικό μέγεθος:} συνολικό πλάτος έως 133 mm, ή 63,5 mm σε στήλη. Τα γράμματα διατηρούν το μέγεθός τους.
\end{itemize}
\BookSection{Β. Άξονας θέσης και μετατόπιση}
\noindent $x_A=-2\,\mathrm{m}$ και $x_B=1\,\mathrm{m}$. Εδώ τονίζεται η μετατόπιση.
\begin{center}
\begin{tikzpicture}[x=1.02cm,y=.95cm]
  \draw[book/axis] (-3.2,0)--(3.5,0) node[book/physlabel,right] {$x\,(\mathrm{m})$};
  \foreach \x in {-3,-2,-1,0,1,2,3}{
    \draw[book/tick] (\x,.07)--(\x,-.07);
    \node[book/ticklabel,below=2pt] at (\x,-.09) {$\x$};
  }
  \node[book/endpoint] at (-2,0) {};
  \node[book/physlabel,above=2pt] at (-2,0) {$A$};
  \node[book/endpoint] at (1,0) {};
  \node[book/physlabel,above=2pt] at (1,0) {$B$};
  \draw[book/mainvector] (-2,.80)--(1,.80)
       node[above,midway,book/physlabel] {$\Delta x=+3\,\mathrm{m}$};
\end{tikzpicture}
\captionof{figure}{Μαύρος άξονας αναφοράς και μετατόπιση με το δεύτερο μελάνι.}
\label{fig:guide-displacement}
\end{center}
\noindent\textbf{Κώδικας:} \texttt{book/axis} για τον άξονα, \texttt{book/mainvector} για το διάνυσμα έμφασης.
\clearpage
\BookSection{Γ. Ελεύθερο διάγραμμα δυνάμεων}
\noindent Το σώμα παριστάνεται ως υλικό σημείο: όλες οι δυνάμεις αρχίζουν στο ίδιο σημείο. Τονίζεται η εφαρμοζόμενη δύναμη $\vec F$. Η ταχύτητα σχεδιάζεται χωριστά, επειδή δεν είναι δύναμη.
\begin{center}
\begin{tikzpicture}[x=1.05cm,y=1cm]
  \draw[book/referencevector] (0,0)--(0,-1.48) node[below,book/physlabel] {$\vec w$};
  \draw[book/mainvector] (0,0)--(2.35,0) node[right,book/physlabel] {$\vec F$};
  \draw[book/referencevector] (0,0)--(-2.4,0) node[left,book/physlabel] {$\vec f$};
  \draw[book/referencevector] (0,0)--(0,1.52) node[above,book/physlabel] {$\vec N$};
  \node[book/endpoint] at (0,0) {};
  \node[book/physlabel] at (.38,-.28) {$m$};
  \draw[book/referencevector] (-.15,2.25)--(1.4,2.25) node[right,book/physlabel] {$\vec\upsilon$};
\end{tikzpicture}
\captionof{figure}{Η δύναμη $\vec F$ είναι το κύριο στοιχείο· οι υπόλοιπες δυνάμεις και η ταχύτητα είναι μαύρα στοιχεία αναφοράς.}
\label{fig:guide-forces}
\end{center}
\BookSection{Δ. Οριζόντια βολή}
\noindent Στο παράδειγμα ορίζονται αρχή στο σημείο βολής, θετικός $x$ δεξιά και θετικός $y$ προς τα κάτω. Τονίζεται η τροχιά.
\begin{center}
\begin{tikzpicture}[x=1.08cm,y=-.69cm]
  \draw[book/axis] (0,0)--(5.0,0) node[above,book/physlabel] {$x$};
  \draw[book/axis] (0,-.16)--(0,3.30) node[left,book/physlabel] {$y$};
  \node[book/endpoint] at (0,0) {};
  \node[book/physlabel,above left] at (0,0) {$O$};
  \draw[book/maincurve,domain=0:4,samples=70,smooth,variable=\t]
     plot ({\t},{.18*\t*\t});
  \draw[book/referencevector] (.05,-.46)--(1.5,-.46)
     node[midway,above,book/physlabel] {$\vec\upsilon_0$};
  \draw[book/guide] (4,0)--(4,2.88);
  \node[book/endpoint] at (4,2.88) {};
  \node[book/physlabel,below] at (4,2.88) {$P$};
\end{tikzpicture}
\captionof{figure}{Τροχιά με δεύτερο μελάνι και αρχική ταχύτητα ως μαύρη αναφορά. Η φορά των αξόνων αφορά το συγκεκριμένο παράδειγμα.}
\label{fig:guide-projectile}
\end{center}
\clearpage
\BookSection{Ε. Γραφικές παραστάσεις}
\noindent Τα TikZ και PGFPlots χρησιμοποιούν τα ίδια πάχη και τις ίδιες αιχμές. Οι αριθμοί των αξόνων είναι 8 pt και οι ετικέτες 9 pt.
\begin{center}
\begin{tikzpicture}
\begin{axis}[book/plot,xmin=0,xmax=8,ymin=0,ymax=6,
  xtick={0,2,4,6,8},ytick={0,1,2,3,4,5,6},
  xlabel={$t\,(\mathrm{s})$},ylabel={$x\,(\mathrm{m})$}]
  \addplot[book/seriesA,domain=0:8,samples=2]{1+0.5*x};
  \addplot[book/seriesPoints] coordinates {(0,1) (4,3) (8,5)};
\end{axis}
\end{tikzpicture}
\captionof{figure}{Θέση–χρόνος: $x=1+0.5t$ σε SI. Η καμπύλη και τα μετρημένα σημεία τονίζονται με δεύτερο μελάνι.}
\label{fig:guide-xt}
\end{center}
\medskip
\begin{center}
\begin{tikzpicture}
\begin{axis}[book/plot,xmin=0,xmax=8,ymin=0,ymax=4,
 xtick={0,2,4,6,8},ytick={0,1,2,3,4},
 xlabel={$t\,(\mathrm{s})$},ylabel={$\upsilon\,(\mathrm{m/s})$}]
 \addplot[book/seriesA,domain=0:8,samples=2]{2};
 \addplot[book/seriesB,domain=0:8,samples=2]{1+0.25*x};
 \node[book/physlabel,above] at (axis cs:6.8,2) {$A$};
 \node[book/physlabel,above] at (axis cs:6.8,2.7) {$B$};
\end{axis}
\end{tikzpicture}
\captionof{figure}{Δύο κινήσεις με το ίδιο μελάνι: συνεχής Α, διακεκομμένη Β. Οι ετικέτες και τα μοτίβα διατηρούν τη διάκριση.}
\label{fig:guide-vt}
\end{center}
\clearpage
\BookSection{ΣΤ. Σχήμα στο πλάτος στήλης}
\noindent Το ακόλουθο διάγραμμα σχεδιάζεται απευθείας σε στήλη 63,5 mm, με γράμματα ίδιου μεγέθους με τα μεγαλύτερα σχήματα.
\par\medskip
\noindent\begin{minipage}[t]{63.5mm}
\vspace{0pt}\centering
\begin{tikzpicture}
\begin{axis}[book/plot-column,xmin=0,xmax=4,ymin=0,ymax=4,
 xtick={0,2,4},ytick={0,2,4},
 xlabel={$t\,(\mathrm{s})$},ylabel={$x\,(\mathrm{m})$}]
 \addplot[book/seriesA,domain=0:4,samples=2]{x};
\end{axis}
\end{tikzpicture}
\captionof{figure}{Γράφημα σε στήλη, με ετικέτες 9 pt και διαβαθμίσεις 8 pt.}
\label{fig:guide-column}
\end{minipage}\hfill
\begin{minipage}[t]{63.5mm}
\vspace{0pt}\small\raggedright
Το πλάτος περιλαμβάνει τις ετικέτες και τις μονάδες. Δεν μικραίνουμε ολόκληρο το σχήμα με \texttt{resizebox}: σχεδιάζουμε στο τελικό του μέγεθος.
\par\smallskip
Χρησιμοποιούμε \texttt{book/plot-column} και λιγότερες διαβαθμίσεις όταν ο διαθέσιμος χώρος είναι μικρός.
\end{minipage}
\par\bigskip
\BookSection{Ζ. Ιδανική τροχαλία}
\noindent Εδώ τονίζονται τα βάρη, ενώ τροχαλία, νήμα και σώματα παραμένουν μαύρα.
\begin{center}
\begin{tikzpicture}[x=1cm,y=1cm]
  \draw[book/body] (-2.0,2.6)--(2.0,2.6);
  \draw[book/referencecurve] (0,2.6)--(0,2.12);
  \draw[book/body] (0,1.6) circle (.52);
  \node[book/endpoint] at (0,1.6) {};
  \draw[book/referencecurve] (-.52,.05)--(-.52,1.6)
    arc[start angle=180,end angle=0,radius=.52]--(.52,-.30);
  \draw[book/body] (-.97,-.65) rectangle (-.07,.05);
  \draw[book/body] (.07,-1.0) rectangle (.97,-.30);
  \node[book/physlabel] at (-.52,-.30) {$m_1$};
  \node[book/physlabel] at (.52,-.65) {$m_2$};
  \draw[book/mainvector] (-.52,-.65)--(-.52,-1.43)
    node[below,book/physlabel] {$\vec w_1$};
  \draw[book/mainvector] (.52,-1.00)--(.52,-1.78)
    node[below,book/physlabel] {$\vec w_2$};
\end{tikzpicture}
\captionof{figure}{Το δεύτερο μελάνι τονίζει τα βάρη σε αυτό το σχήμα· δεν αποτελεί σταθερό χρώμα του φυσικού μεγέθους.}
\label{fig:guide-pulley}
\end{center}
\clearpage
\BookSection{Η. Εντολές και έλεγχος}
\begin{BookPlainTable}
\begin{tabularx}{\linewidth}{@{}p{29mm}>{\raggedright\arraybackslash}X@{}}
Άξονες & \texttt{book/axis}, \texttt{book/tick} \\
Διανύσματα & \texttt{book/mainvector}, \texttt{book/secondaryvector}, \texttt{book/referencevector} \\
Καμπύλες & \texttt{book/maincurve}, \texttt{book/secondarycurve}, \texttt{book/referencecurve} \\
Σημεία & \texttt{book/endpoint}, \texttt{book/pointMain} \\
Ετικέτες & \texttt{book/physlabel}, \texttt{book/ticklabel} \\
Βοηθητικά & \texttt{book/guide} \\
Γραφήματα & \texttt{book/plot}, \texttt{book/plot-column} \\
Σειρές & \texttt{book/seriesA}, \texttt{book/seriesB} \\
\end{tabularx}
\end{BookPlainTable}
\par\medskip
\noindent Κάθε σχήμα έχει λεζάντα, \texttt{label} και αναφορά από το κείμενο.
\par\medskip
Πριν από την παράδοση ελέγχουμε το τελικό μέγεθος, τις μονάδες SI, τις θετικές φορές, τα σημεία εφαρμογής των δυνάμεων και τη διάκριση γραμμών. Η PRINT πρέπει να διατηρεί μόνο τις δύο προβλεπόμενες πλάκες.
